\chapter{Sonuçlar ve Doğrulama}
\label{chapter:results}

Bu bölümde, geliştirilen Ay yörünge simülasyonu üzerinde doğrulama ve duyarlılık testleri gerçekleştirilmiştir. Ele alınan başlıca konular şunlardır: (i) Ay sferik harmonik (SH) yerçekimi modelinde derece kesiminin (truncation) yörünge dinamiğine etkisi ve mascon kaynaklı bozucu etkilerin hangi çözünürlükte belirginleştiği, (ii) SRP, Albedo ve Termal IR gibi ışınım kökenli pertürbasyonların birikimli sürüklenme (drift) etkileri, (iii) farklı integratörlerin aynı başlangıç koşullarında doğruluk--hesap maliyeti açısından karşılaştırılması ve (iv) tutulma (eclipse) süresi ile belirli koordinatlardan geçiş sayısı gibi operasyonel metriklerin değerlendirilmesi.

\newpage
\section{Ay Yerçekimi Modeli: Harmonik Derecenin Etkisi}
\label{sec:gravity_degree_effect}

\subsection{Kısa Süreli Duyarlılık: 6 Saatlik Yayılım}
Aynı başlangıç koşulları altında, farklı SH derece kesimleri ($N=\{0,2,10,50,180\}$) kullanılarak 6 saatlik yayılımlar yapılmış; irtifa evrimi ve yer izi (ground track) davranışı karşılaştırılmıştır. Şekil~\ref{fig:truncation_effect}’te görüldüğü üzere, yalnızca merkezi terimin ($N=0$) kullanıldığı durumda yörünge dinamiği idealize olmakta ve irtifa değişimi ihmal edilebilir düzeyde kalmaktadır. Derece arttıkça masconların temsil çözünürlüğü yükselmekte; buna bağlı olarak kısa zaman ölçeğinde dahi irtifa salınımlarının genliği ve fazı belirgin şekilde değişmektedir. Yüksek derecelerde eğrilerin birbirine yaklaşması kısa süreli yayılımlarda sayısal yakınsama eğilimine işaret etse de, uzun süreli propagasyonlarda bu küçük farkların birikerek anlamlı konum/irtifa sapmalarına yol açtığı gözlenmiştir.

\begin{figure}[htbp]
    \centering
    \includegraphics[width=\linewidth]{Resimler/Results/Figure_1.png}
    \caption{Sferik harmonik derece kesiminin (truncation) yörünge irtifa evrimi ve yer izi üzerindeki etkisi.
    Sol: farklı $N$ değerleri için irtifa değişimi. Sağ: $N=180$ için sabit referans çerçevesinde yer izi (renk: zaman).}
    \label{fig:truncation_effect}
\end{figure}

\subsection{Uzun Süreli Yakınsama: Yörünge Ömrü}
\label{subsec:convergence}
Yerçekimi model derecesinin görev ömrü tahminine etkisini nicel olarak değerlendirmek amacıyla, mascon etkilerine duyarlı bir senaryo seçilmiştir: $i=90^\circ$ (polar), başlangıç irtifası $h_0=50$ km olan dairesel yörünge. Farklı $N$ değerleriyle yürütülen propagasyonlarda çarpışma eşiğine (impact threshold) ulaşma zamanı “yörünge ömrü” olarak tanımlanmıştır. Tablo~\ref{tab:convergence}’de görüldüğü üzere, düşük dereceli modeller ($N<50$) mascon kaynaklı yüksek frekanslı yerçekimi bileşenlerini yeterince temsil edemediğinden yörünge ömrünü iyimser tahmin etme eğilimindedir. $N \approx 150$--$180$ aralığında sonuçların stabilize olması, bu senaryo için derecesel yakınsamaya ulaşıldığını göstermektedir.

\begin{table}[htbp]
\centering
\small
\caption{Sferik harmonik derecesine ($N$) göre yörünge ömrü yakınsaması (polar, $h_0=50$ km).}
\label{tab:convergence}
\begin{tabular}{c c p{7.4cm}}
\toprule
\textbf{Model Derecesi ($N$)} & \textbf{Yörünge Ömrü (saat)} & \textbf{Gözlem ve teknik yorum} \\
\midrule
10--20 & $\sim$250+ & Alan aşırı homojen temsil edilir; mascon etkileri bastırılır ve bozunma hızı olduğundan düşük çıkar. \\
\addlinespace
50 & 210.5 & Bölgesel anomaliler belirginleşir; periselenede düşüş hızlanır. \\
\addlinespace
120 & 185.2 & Yüksek frekanslı bileşenler devreye girer; irtifa kayıpları daha keskin hale gelir. \\
\addlinespace
150--180 & \textbf{180.4} & \textbf{Yakınsama bölgesi:} $N>180$ için sonuç değişimi sınırlıdır; model doyuma ulaşır. \\
\bottomrule
\end{tabular}
\end{table}

\subsection{Fiziksel Mekanizma: Bozunma ve Çarpışmaya Gidiş}
Simülasyon çıktıları, 50 km irtifadaki polar yörüngenin kararlılığını yitirmesinde aşağıdaki mekanizmaların baskın olduğunu göstermektedir:
\begin{itemize}
    \item \textbf{Eksantrisite büyümesi:} Mascon kaynaklı düzensiz yerçekimi alanı, her turda yörünge elemanlarında enerji ve açısal momentum alışverişine neden olarak dairesellikten sapmayı artırır; periselene irtifası kademeli olarak düşer.
    \item \textbf{Yüksek frekanslı gradyan etkileri:} $N=180$ gibi yüksek dereceler, keskin yerçekimi gradyanlarını içerdiğinden bazı bölgelerde ani irtifa kayıpları ve faz kaymaları tetiklenebilir.
    \item \textbf{Kritik eşik ve son faz:} Periselene irtifası belirlenen impact threshold seviyesine indiğinde, temas/çarpışma süreci hızla gerçekleşir ve senaryoda toplam görev ömrü yaklaşık 180.4 saat olarak bulunur.
\end{itemize}

\subsection{Özet}
Bu analiz, düşük çözünürlüklü gravite modellerinin ($N<50$) yörünge ömrü tahminlerinde anlamlı düzeyde iyimser hatalara yol açabildiğini; görev planlama aşamasında derece yakınsamasının mutlaka test edilmesi gerektiğini göstermektedir. İncelenen polar 50 km senaryosu için güvenli görev penceresi yaklaşık 180 saat mertebesindedir.


%===========================================================
% Ay Küresel Harmonik Yerçekimi için
% Hibrit Truncation (Kesme) Hata Tahmini
%===========================================================

\section{Yerçekimi Modeli Truncation Hatası ve Hibrit Kuyruk Ekstrapolasyonu}

\subsection{Küresel harmonik yerçekimi modeli ve kesme}
Tam normalize edilmiş küresel harmoniklerle ifade edilen bir yerçekimi potansiyeli şu şekilde yazılabilir:
\begin{equation}
V(r,\phi,\lambda)
=
\frac{\mu}{r}
\left[
1
+
\sum_{n=2}^{N}
\left(\frac{R}{r}\right)^{n}
\sum_{m=0}^{n}
\bar{P}_{nm}(\sin\phi)
\left(
\bar{C}_{nm}\cos m\lambda
+
\bar{S}_{nm}\sin m\lambda
\right)
\right],
\label{eq:V_sh}
\end{equation}
burada $\mu=GM$ yerçekimi parametresi, $R$ referans yarıçapı, $r$ merkezden uzaklık, $\phi$ enlem, $\lambda$ boylam ve $\bar{P}_{nm}$ tam normalize edilmiş Legendre fonksiyonlarıdır. Modelde kullanılan en yüksek derece $N$ ile gösterilir.

Pratikte, mevcut yerçekimi modeli sonlu bir derece $N_{\mathrm{cut}}$’ta kesilir. Bu durumda $n > N_{\mathrm{cut}}$ olan terimlerin katkısı \emph{kesme (omission) hatası} olarak adlandırılır.

\subsection{Noktasal “gerçek” truncation hatası (yön-bağımlı en kötü durum)}
Mevcut en yüksek dereceli çözüm $N_{\mathrm{truth}}$ (bu çalışmada $N_{\mathrm{truth}}=180$) referans alınarak, herhangi bir $\mathbf{r}$ konumundaki ivme farkı
\begin{equation}
\Delta\mathbf{a}(\mathbf{r};N_{\mathrm{cut}})
=
\mathbf{a}(\mathbf{r};N_{\mathrm{cut}})
-
\mathbf{a}(\mathbf{r};N_{\mathrm{truth}})
\label{eq:delta_a}
\end{equation}
şeklinde tanımlanır.

Buna karşılık bağıl ivme hatası metriği
\begin{equation}
\varepsilon(\mathbf{r};N_{\mathrm{cut}})
=
\frac{\left\|\Delta\mathbf{a}(\mathbf{r};N_{\mathrm{cut}})\right\|}
{\left\|\mathbf{a}(\mathbf{r};N_{\mathrm{truth}})\right\|}
\label{eq:rel_err_pointwise}
\end{equation}
olarak verilir.

Truncation hatası konuma (enlem/boylam) bağlı olduğundan, küre üzerinde yaklaşık eşit dağılmış $N_{\mathrm{dir}}$ adet yön vektörü $\{\hat{\mathbf{u}}_k\}$ (örneğin Fibonacci küresi) kullanılır. Yüzeyden $h$ irtifasında ($r=R+h$) en kötü durum hatası şu şekilde tanımlanır:
\begin{equation}
\varepsilon_{\max}(h;N_{\mathrm{cut}})
=
\max_{k}
\varepsilon\!\left(r\,\hat{\mathbf{u}}_k;N_{\mathrm{cut}}\right).
\label{eq:worst_case_err}
\end{equation}

Belirli bir hata eşiği $\tau$ için gerekli minimum kesme derecesi
\begin{equation}
N_{\max}(h;\tau)
=
\min\left\{
N_{\mathrm{cut}} \;:\;
\varepsilon_{\max}(h;N_{\mathrm{cut}})\le \tau
\right\}
\label{eq:Nmax_def}
\end{equation}
olarak tanımlanır.

\paragraph{Not (metrik uyumsuzluğu).}
$\varepsilon_{\max}$ metriği \emph{noktasal ve en kötü durum} yaklaşımıdır. Buna karşılık Kaula tabanlı modeller istatistiksel (RMS) temellidir. Bu nedenle Kaula parametrelerinin doğrudan $\varepsilon_{\max}$ ile kalibre edilmesi matematiksel olarak tutarlı değildir.

\subsection{Derece güç spektrumu ve ampirik katsayı RMS’i}
Bir derece için güç spektrumu şu şekilde tanımlanır:
\begin{equation}
P_n
=
\sum_{m=0}^{n}
\left(\bar{C}_{nm}^2+\bar{S}_{nm}^2\right).
\label{eq:power_spectrum}
\end{equation}

Buna karşılık her bir katsayı için ampirik RMS büyüklüğü
\begin{equation}
\sigma_{\mathrm{coef,emp}}(n)
=
\sqrt{\frac{P_n}{2n+1}}
\label{eq:sigma_emp}
\end{equation}
olarak elde edilir.

\subsection{Kaula tipi güç yasası ve üst derece devamı}
Kaula yaklaşımında katsayı RMS’lerinin derece ile şu şekilde azaldığı varsayılır:
\begin{equation}
\sigma_{\mathrm{coef}}(n)
\approx
\frac{A}{n^{p}}.
\label{eq:kaula_law}
\end{equation}

Bu çalışmada $A$ ve $p$ sabitleri literatürden alınmak yerine, mevcut çözümün ($n \le 180$) üst bantındaki spektrumdan elde edilir. Logaritmik ölçekte
\begin{equation}
\log \sigma_{\mathrm{coef,emp}}(n)
\approx
\log A - p\,\log n
\label{eq:loglog_fit}
\end{equation}
bağıntısı doğrusal regresyonla fit edilir.

\subsection{Hibrit RMS truncation hata modeli}
Mevcut en yüksek derece $N_{\mathrm{truth}}=180$ ve ekstrapolasyon üst sınırı $N_{\mathrm{total}}=1200$ olmak üzere, derece genliği
\begin{equation}
\sigma_{\mathrm{deg}}(n)
=
\begin{cases}
\sqrt{P_n}, & 2\le n \le N_{\mathrm{truth}},\\[6pt]
\sqrt{2n+1}\,\dfrac{A}{n^{p}}, & N_{\mathrm{truth}}< n \le N_{\mathrm{total}},
\end{cases}
\label{eq:sigma_deg_hybrid}
\end{equation}
şeklinde tanımlanır.

Her derecenin ivmeye katkısı yaklaşık olarak
\begin{equation}
t_n(r)
=
(n+1)\,\sigma_{\mathrm{deg}}(n)
\left(\frac{R}{r}\right)^{n}
\label{eq:tn_model}
\end{equation}
ile modellenir.

Buna göre RMS truncation hatası
\begin{equation}
\varepsilon_{\mathrm{RMS}}(r;N_{\mathrm{cut}})
=
\sqrt{
\sum_{n=N_{\mathrm{cut}}+1}^{N_{\mathrm{total}}}
t_n(r)^2
}
\label{eq:rms_tail_error}
\end{equation}
olarak hesaplanır.

\subsection{Yorum ve kullanım önerisi}
Bu yaklaşım iki tamamlayıcı çıktı üretir:
\begin{itemize}
\item \textbf{Gerçek (en kötü durum) truncation hatası:} $N \le 180$ için doğrudan model farkından hesaplanır.
\item \textbf{Hibrit RMS truncation hatası:} 180’e kadar gerçek spektrum, daha üst dereceler için Kaula-tabanlı devam kullanılır.
\end{itemize}

Özellikle iniş ve alçak irtifa operasyonlarında yüksek dereceli terimler baskın hale geldiğinden, hibrit RMS model gerekli minimum derece seçimi için fiziksel ve sayısal olarak tutarlı bir yöntem sunar. Emniyet-kritik uygulamalarda RMS değerlerin üzerine ilave muhafazakâr çarpanlar eklenmesi önerilir.

%===========================================================
% Son
%===========================================================

\begin{figure}[H]
    \centering
    \includegraphics[width=\textwidth]{Resimler/Results/true_truncation_error_l180.png}
    \caption{%
    Ay yerçekimi modeli için gerçek truncation hatasının,
    irtifa ve küresel harmonik kesme derecesine bağlı olarak
    değişimi.
    Renk haritası, $N_{\mathrm{truth}}=180$ dereceli çözüm referans
    alınarak hesaplanan maksimum bağıl ivme hatasını
    ($\max \|\Delta\mathbf{a}\|/\|\mathbf{a}\|$) göstermektedir.
    Kesikli çizgiler, farklı hata eşiklerine
    ($10^{-6}$, $10^{-9}$, $10^{-12}$, $10^{-15}$) karşılık gelen
    gerekli minimum kesme derecelerini belirtmektedir.
    }
    \label{fig:true_truncation_error}
\end{figure}


\begin{figure}[H]
    \centering
    \includegraphics[width=0.9\textwidth]{Resimler/Results/hybrid_rms_tail_error_to1200.png}
    \caption{%
    Belirlenen bağıl ivme hata eşiklerine göre,
    irtifaya bağlı olarak gerekli minimum küresel harmonik
    kesme derecesi $N_{\max}$.
    Eğriler, en kötü durum (worst-case) hata metriği
    kullanılarak elde edilmiştir.
    }
    \label{fig:nmax_vs_altitude}
\end{figure}

\begin{table}[H]
\centering
\caption{%
Gerçek (model tabanlı) yerçekimi çözümü kullanılarak,
farklı irtifalarda belirlenen bağıl ivme hata eşiklerini
sağlamak için gerekli minimum küresel harmonik kesme
dereceleri.
Hata metriği, tüm uzaysal yönler arasında maksimum
(worst-case) bağıl ivme hatası olarak tanımlanmıştır.
}
\label{tab:true_required_degree}
\begin{tabular}{c|ccccccc}
\hline
\textbf{İrtifa (km)} &
\textbf{$10^{-3}$} &
\textbf{$10^{-5}$} &
\textbf{$10^{-6}$} &
\textbf{$10^{-7}$} &
\textbf{$10^{-9}$} &
\textbf{$10^{-12}$} &
\textbf{$10^{-15}$} \\
\hline
10    & 50  & 180 & 180 & 180 & 180 & 180 & 180 \\
50    & 9   & 139 & 179 & 180 & 180 & 180 & 180 \\
100   & 3   & 72  & 105 & 144 & 180 & 180 & 180 \\
200   & 2   & 38  & 57  & 77  & 116 & 176 & 180 \\
500   & 2   & 13  & 23  & 33  & 50  & 77  & 104 \\
1000  & 2   & 8   & 12  & 17  & 27  & 42  & 57 \\
2000  & 2   & 4   & 7   & 10  & 16  & 25  & 34 \\
\hline
\end{tabular}
\end{table}



\begin{table}[H]
\centering
\caption{%
Hibrit truncation hata modeli kullanılarak,
$n \le 180$ için gerçek yerçekimi spektrumu,
$n > 180$ için Kaula-tabanlı spektral devam
(1200. dereceye kadar) ile hesaplanan
gerekli minimum küresel harmonik kesme
dereceleri.
Hata metriği RMS bağıl ivme hatasıdır.
}
\label{tab:hybrid_required_degree}
\begin{tabular}{c|ccccccc}
\hline
\textbf{İrtifa (km)} &
\textbf{$10^{-3}$} &
\textbf{$10^{-5}$} &
\textbf{$10^{-6}$} &
\textbf{$10^{-7}$} &
\textbf{$10^{-9}$} &
\textbf{$10^{-12}$} &
\textbf{$10^{-15}$} \\
\hline
10    & 2   & 425 & 747  & 1077 & 1200 & 1200 & 1200 \\
50    & 2   & 102 & 170  & 241  & 389  & 618  & 852  \\
100   & 2   & 53  & 89   & 126  & 202  & 318  & 437  \\
200   & 2   & 26  & 46   & 66   & 106  & 166  & 227  \\
500   & 2   & 10  & 19   & 28   & 45   & 72   & 99   \\
1000  & 2   & 5   & 10   & 15   & 25   & 40   & 55   \\
2000  & 2   & 3   & 6    & 9    & 14   & 24   & 32   \\
4000  & 2   & 2   & 3    & 5    & 9    & 15   & 21   \\
\hline
\end{tabular}
\end{table}




Tablo~\ref{tab:true_required_degree} ve
Tablo~\ref{tab:hybrid_required_degree} birlikte
incelendiğinde, truncation hatasının hem
irtifaya hem de kullanılan hata metriğine
son derece duyarlı olduğu görülmektedir.
Gerçek model tabanlı analiz, worst-case
yaklaşımı nedeniyle daha konservatif
sonuçlar üretirken; hibrit model RMS
metrik kullanarak istatistiksel ortalama
davranışı temsil etmektedir.
Özellikle alçak irtifalarda,
mevcut 180 dereceli yerçekimi çözümlerinin
iniş ve yüzeye yakın uçuş senaryoları için
yetersiz olduğu, yüksek dereceli harmonik
içeriğin ise Kaula-tabanlı spektral devam
yöntemleri ile tahmin edilmesinin
kaçınılmaz olduğu ortaya konmuştur.
